\documentclass[12pt]{article}
\usepackage[utf8]{inputenc}
\usepackage[T1]{fontenc}
\usepackage{hyperref}
\title{Yeri Language Resources for K-12 Curriculum Development}
\author{Compiled by Claude Sonnet 5.5 (claude-sonnet-5-5) for Andr\'as Kornai}
\date{}
\begin{document}
\maketitle

\section{Introduction}
Yeri (also Yapunda, Reiwo, Hiro; ISO 639-3 \texttt{yev}, Glottocode yapu1240) is a severely endangered Torricelli language of Sandaun Province, Papua New Guinea. It is tied to a single village, Yapunda. Estimates of speakers range from 10--99 (UNESCO) to 100--150 (Wikipedia). A 2014 personal communication cited by the Endangered Languages Project says ``probably closer to 60''. Fluent speakers are mostly over 40, and children speak Tok Pisin.

\section{Primary Educational Resources}
No primers, textbooks or school programmes were found in the search results.

\section{Linguistic and Reference Materials}
\begin{itemize}
\item A dissertation grammar based on eleven months of fieldwork, described as the first description of Yeri. It notes rare verb infixation after the first syllable of the stem: \url{https://www.proquest.com/docview/1878004238}
\item Max Planck Institute for Evolutionary Anthropology, ``A Grammar of Yeri'' project page: \url{https://www.eva.mpg.de/linguistics/past-research-resources/documentation-and-description/a-grammar-of-yeri/}
\item Wikipedia, Yapunda language: \url{https://en.wikipedia.org/wiki/Yapunda_language}
\end{itemize}

\section{Religious and Community Resources}
\begin{itemize}
\item Joshua Project, Yapunda in Papua New Guinea: \url{https://joshuaproject.net/people_groups/15990}
\end{itemize}
No scripture translation or community organisation was found.

\section{Digital and Audio Resources}
The search results showed no audio or corpus holdings. The Endangered Languages Project entry may link to resources, but this was not confirmed:
\begin{itemize}
\item \url{https://endangeredlanguages.com/lang/2709}
\end{itemize}

\section{Curriculum Development Context}
The language is not known to be taught in schools, and only older adults use it. The scholarly grammar is the main resource available. Any curriculum would be a revitalisation effort built with the Yapunda community, most likely as oral and cultural materials alongside Tok Pisin. A writing system and teaching materials are not documented in what was found.
\end{document}